\chapter{Contexto y objetivos}

El trabajo se desarrolla como TFG del Grado en Ingeniería Informática, mención de Ingeniería del Software. Tras la modificación aprobada el 19 de junio de 2026, el objetivo dejó de ser una plataforma de infraestructura cloud para centrarse en el diseño y desarrollo integral de la aplicación. La implementación anterior se conserva como fuente de conocimiento, pero el producto activo se reconstruye desde cero para que cada incremento recorra realmente todas las capas.

\section{Objetivo general}

Diseñar, construir, probar y desplegar una aplicación web mantenible que permita gestionar las operaciones esenciales de una empresa de transporte y que constituya una demostración verificable del ciclo de vida completo del software.

\section{Objetivos específicos}

\begin{itemize}
  \item Obtener y formalizar requisitos trazables desde una entrevista con la parte interesada.
  \item Diseñar el dominio, las interfaces y una arquitectura separada de backend y frontend.
  \item Implementar incrementos verticales utilizables, con autenticación y autorización desde el comienzo.
  \item Automatizar pruebas de comportamiento en backend y frontend.
  \item Garantizar ejecución local reproducible mediante contenedores e integración continua.
  \item Desplegar el producto completo en un entorno accesible y evaluar los resultados.
\end{itemize}

El éxito se medirá por el cumplimiento de RF-01 a RF-14, la ejecución satisfactoria de los cuatro flujos de negocio, la ausencia de secretos en el repositorio y la posibilidad de reproducir build, pruebas y despliegue siguiendo documentación explícita.
